% Part: normal-modal-logic
% Chapter: frame-definability
% Section: first-order-definability

\documentclass[../../../include/open-logic-section]{subfiles}

\begin{document}

\olfileid{nml}{frd}{fol}

\olsection{قابلية التعريف من الرتبة الأولى}

سبق أن عرّفنا جملة من خصائص الوصول بـ!!{formula}s موجهية؛
فالتناظر، مثلًا، تعرّفه ال!!{formula}~\Ax{B}، أي~$\OLId{latin-ordinary}{p} \lif \Box\Diamond \OLId{latin-ordinary}{p}$.
لكن الشروط السابقة كلها تقبل أيضًا التعبير بـ!!{formula}s من
الرتبة الأولى، في !!a{language} لا تضم إلا !!{predicate} ثنائيًا
واحدًا. فقولنا
$\lforall[\OLId{latin-ordinary}{x}][\lforall[\OLId{latin-ordinary}{y}][(\Atom{Q}{x,y} \lif \Atom{Q}{y, x})]]$
يعرّف التناظر: ذلك أن !!{structure} من الرتبة الأولى~$\Struct{M}$،
حيث~$\Domain{\OLId{latin-looped}{M}} = \OLId{latin-looped}{W}$ و$\Assign{\OLId{latin-looped}{Q}}{\OLId{latin-looped}{M}} = \OLId{latin-looped}{R}$، تحقق هذه الصيغة
متى وفقط متى كانت~$\OLId{latin-looped}{R}$ متناظرة. فنضع لهذا المعنى التعريف الآتي:

\begin{defn}
  يكون صنف الأطر~$\mClass{F}$ \emph{قابلًا للتعريف من الرتبة الأولى}
  إذا وُجدت !!a{sentence}~$!A$ في لغة الرتبة الأولى ذات
  !!{predicate} ثنائي واحد~$\OLId{latin-looped}{Q}$، بحيث
  $\mModel{F} = \tuple{\OLId{latin-looped}{W}, \OLId{latin-looped}{R}} \in \mClass{F}$ إذا وفقط إذا كان
  $\Sat{\OLId{latin-looped}{M}}{!A}$ في ال!!{structure} من الرتبة الأولى~$\Struct{M}$ ذات
  $\Domain{\OLId{latin-looped}{M}} = \OLId{latin-looped}{W}$ و$\Assign{\OLId{latin-looped}{Q}}{\OLId{latin-looped}{M}} = \OLId{latin-looped}{R}$.
\end{defn}

غير أن اتفاق الطريقتين في الخصائص وال!!{formula}s الموجهية
السابقة لا يعمّ جميع الأحوال. فقد تعرّف !!{formula} موجهية صنفًا
لا يُعرّف من الرتبة الأولى؛ وقد يكون صنف من الأطر قابلًا للتعريف
من الرتبة الأولى، ولا يكون قابلًا للتعريف بصيغة موجهية.

يعطي مثالًا مضادًا للدعوى الأولى صيغة لوب:
\begin{equation}
\Box(\Box \OLId{latin-ordinary}{p} \lif \OLId{latin-ordinary}{p}) \lif \Box \OLId{latin-ordinary}{p}. \tag{\Ax{W}}
\end{equation}
تعرّف~\Ax{W} صنف الأطر المتعدية والمعكوسة الحسنة التأسيس. وتكون
العلاقة حسنة التأسيس إذا لم توجد متتالية لا نهائية
$\OLId{latin-ordinary}{w}_\OLMathNumeral{1}$، و$\OLId{latin-ordinary}{w}_\OLMathNumeral{2}$، و\dots{} بحيث~$Rw_\OLMathNumeral{2}\OLId{latin-ordinary}{w}_\OLMathNumeral{1}$، و$Rw_\OLMathNumeral{3}\OLId{latin-ordinary}{w}_\OLMathNumeral{2}$، و\dots. فمثلًا،
العلاقة~$<$ على~$\Nat$ حسنة التأسيس، بينما العلاقة~$<$ على~$\Int$
ليست كذلك. وتكون العلاقة معكوسة حسنة التأسيس إذا وفقط إذا كان
معكوسها حسن التأسيس. ومن ثم فالعلاقات المعكوسة الحسنة التأسيس هي التي
لا توجد فيها متتالية لا نهائية~$\OLId{latin-ordinary}{w}_\OLMathNumeral{1}$، و$\OLId{latin-ordinary}{w}_\OLMathNumeral{2}$، و\dots{} بحيث
$Rw_\OLMathNumeral{1}\OLId{latin-ordinary}{w}_\OLMathNumeral{2}$، و$Rw_\OLMathNumeral{2}\OLId{latin-ordinary}{w}_\OLMathNumeral{3}$، و\dots.

لكن لا توجد !!{formula} من الرتبة الأولى تعرّف العلاقات المتعدية
المعكوسة الحسنة التأسيس. افترض، خلاف ذلك، أن~$\Sat{\OLId{latin-looped}{M}}{!F}$ إذا وفقط
إذا كانت~$\OLId{latin-looped}{R} = \Assign{\OLId{latin-looped}{Q}}{\OLId{latin-looped}{M}}$ متعدية ومعكوسة حسنة التأسيس. ولكل
$\OLId{latin-ordinary}{n} \ge \OLMathNumeral{2}$، لتكن~$!A_\OLId{latin-ordinary}{n}$ ال!!{formula}
\[
(\Atom{Q}{a_1,a_2} \land \dots \land
    \Atom{Q}{a_{n-1},a_{n}})
\]
وتأمل الآن مجموعة ال!!{formula}s
\[
\OLId{greek-variable}{Gamma} = \{!F, !A_\OLMathNumeral{2}, !A_\OLMathNumeral{3}, \dots\}.
\]
كل مجموعة جزئية منتهية~$\OLId{greek-variable}{Delta}$ من~$\OLId{greek-variable}{Gamma}$ قابلة للتحقيق. نأخذ
$\OLId{latin-ordinary}{k} \ge \OLMathNumeral{2}$ أكبر عدد وردت~$!A_\OLId{latin-ordinary}{k}$ له في~$\OLId{greek-variable}{Delta}$؛ وإن لم ترد فيها
صيغة من هذا النوع، أخذنا~$\OLId{latin-ordinary}{k}=\OLMathNumeral{2}$. ولنجعل
$\Domain{\OLId{latin-looped}{M}_\OLId{latin-ordinary}{k}} = \{\OLMathNumeral{1}, \dots, \OLId{latin-ordinary}{k}\}$، و$\Assign{\OLId{latin-ordinary}{a}_\OLId{latin-ordinary}{i}}{\OLId{latin-looped}{M}_\OLId{latin-ordinary}{k}} = \OLId{latin-ordinary}{i}$
للثوابت التي تحتاجها~$\OLId{greek-variable}{Delta}$، و$\Assign{\OLId{latin-looped}{Q}}{\OLId{latin-looped}{M}_\OLId{latin-ordinary}{k}} = <$؛
أما بقية الثوابت فتعين اعتباطًا في المجال.
ولما كانت~$<$ على~$\{\OLMathNumeral{1}, \dots, \OLId{latin-ordinary}{k}\}$ متعدية ومعكوسة حسنة التأسيس،
حصل~$\Sat{\OLId{latin-looped}{M}_\OLId{latin-ordinary}{k}}{!F}$. ويعطي البناء~$\Sat{\OLId{latin-looped}{M}_\OLId{latin-ordinary}{k}}{!A_\OLId{latin-ordinary}{i}}$
لكل~$!A_\OLId{latin-ordinary}{i} \in \OLId{greek-variable}{Delta}$. فمن تراص منطق الرتبة الأولى تتحقق~$\OLId{greek-variable}{Gamma}$
في !!{structure} ما~$\Struct{M}$. وحينئذ يلزم من~$\Sat{\OLId{latin-looped}{M}}{!F}$
أن تكون~$\Assign{\OLId{latin-looped}{Q}}{\OLId{latin-looped}{M}}$ معكوسة حسنة التأسيس.
غير أن~$\Assign{\OLId{latin-ordinary}{a}_\OLMathNumeral{1}}{\OLId{latin-looped}{M}}$، و$\Assign{\OLId{latin-ordinary}{a}_\OLMathNumeral{2}}{\OLId{latin-looped}{M}}$، و\dots{}
سلسلة لا نهائية من النوع الممنوع، وهذا تناقض.

وأما المثال المضاد في الاتجاه الآخر فخاصية الكلية: لكل~$\OLId{latin-ordinary}{u}$ و$\OLId{latin-ordinary}{v}$
يكون~$Ruv$. وهي من الرتبة الأولى، إذ تعرّفها
$\lforall[\OLId{latin-ordinary}{x}][\lforall[\OLId{latin-ordinary}{y}][\Atom{Q}{x,y}]]$. ومع ذلك لا توجد صيغة
موجهية تصلح في الأطر الكلية وحدها جميعًا. والسبب نتيجة لها
أهميتها: ما يصلح في جميع الأطر الكلية هو بعينه ما يصلح في جميع
الأطر الانعكاسية المتناظرة المتعدية. ولما كان في هذه الأخيرة أطر
غير كلية، كانت كل~!!{formula} صالحة في الأطر الكلية صالحة أيضًا
في بعض الأطر غير الكلية، فلا تميز الكلية.

\end{document}
